\chapter{Módulo 8 --- Reportes reproducibles con Quarto y proyecto
integrador}

\begin{itemize}
\tightlist
\item
  \textbf{Duración estimada:} 1 h teoría + 4 h proyecto
\item
  \textbf{Objetivos de aprendizaje.} Al finalizar, la persona podrá:

  \begin{enumerate}
  \def\labelenumi{\arabic{enumi}.}
  \tightlist
  \item
    \textbf{Crear} un documento Quarto (\passthrough{\lstinline!.qmd!})
    que combine texto, código R y resultados.
  \item
    \textbf{Configurar} opciones de \emph{chunks}
    (\passthrough{\lstinline!echo!}, \passthrough{\lstinline!warning!},
    \passthrough{\lstinline!fig-cap!}) y formatos de salida (HTML, Word,
    PDF).
  \item
    \textbf{Organizar} un proyecto reproducible (estructura de carpetas,
    rutas relativas, \passthrough{\lstinline!renv!}).
  \item
    \textbf{Diseñar} y \textbf{producir} un análisis completo de
    principio a fin (proyecto integrador).
  \end{enumerate}
\end{itemize}

\section{Contenidos}

\begin{enumerate}
\def\labelenumi{\arabic{enumi}.}
\tightlist
\item
  Por qué la reproducibilidad: el código \emph{es} la documentación del
  análisis.
\item
  Anatomía de un \passthrough{\lstinline!.qmd!}: encabezado YAML,
  Markdown y \emph{chunks} de código.
\item
  Opciones de \emph{chunk} con \passthrough{\lstinline!\#|!}.
\item
  Tablas con \passthrough{\lstinline!knitr::kable()!}; gráficos con
  leyenda.
\item
  Renderizar: botón \textbf{Render} en RStudio o
  \passthrough{\lstinline!quarto render!} en terminal.
\item
  Estructura de proyecto y control de dependencias con
  \passthrough{\lstinline!renv!}.
\item
  Guía de estilo del tidyverse; control de versiones con Git
  (introducción).
\end{enumerate}

\section{Desarrollo}

\subsection{Un documento Quarto mínimo}

Guarda esto como \passthrough{\lstinline!reporte.qmd!} dentro del
proyecto y presiona \textbf{Render}:

\begin{lstlisting}
---
title: "Reporte de ventas 2025"
author: "Tu nombre"
format:
  html:
    toc: true
execute:
  warning: false
---

## Resumen

```{r}
#| label: datos
#| echo: false
library(tidyverse)
ventas <- read_csv("datos/ventas.csv", show_col_types = FALSE) |>
  mutate(ingreso = precio * unidades)
total <- sum(ventas$ingreso, na.rm = TRUE)
```

El ingreso total del año fue de **`r scales::dollar(total)`**.

```{r}
#| label: fig-producto
#| fig-cap: "Ingreso anual por producto"
ventas |>
  summarise(ingreso = sum(ingreso, na.rm = TRUE), .by = producto) |>
  ggplot(aes(ingreso, fct_reorder(producto, ingreso))) +
  geom_col() +
  labs(x = "Ingreso", y = NULL)
```
\end{lstlisting}

\begin{itemize}
\tightlist
\item
  \passthrough{\lstinline!`r expresion`!} inserta resultados
  \textbf{dentro del texto}: si cambian los datos, el número se
  actualiza solo.
\item
  Cambia \passthrough{\lstinline!format: html!} por
  \passthrough{\lstinline!docx!} o \passthrough{\lstinline!pdf!} (este
  último requiere TinyTeX:
  \passthrough{\lstinline!quarto install tinytex!}).
\end{itemize}

\subsection{Proyecto reproducible}

\begin{lstlisting}
mi-analisis/
|-- mi-analisis.Rproj
|-- datos/          # datos crudos: nunca se editan a mano
|-- R/              # funciones reutilizables
|-- reporte.qmd
`-- salidas/        # gráficos y tablas generadas
\end{lstlisting}

\begin{lstlisting}[language=R]
# Registrar versiones de paquetes del proyecto (una vez, en la consola):
# install.packages("renv"); renv::init()
# Después de instalar o actualizar paquetes:
# renv::snapshot()
\end{lstlisting}

\section{Proyecto integrador (capstone)}

\textbf{Caso:} Eres analista de una cadena de electrónica. Dirección
pide un \textbf{reporte de desempeño 2025} con recomendaciones para la
temporada alta de 2026.

\textbf{Entregables:}

\begin{enumerate}
\def\labelenumi{\arabic{enumi}.}
\tightlist
\item
  Proyecto de RStudio con la estructura anterior.
\item
  \passthrough{\lstinline!reporte.qmd!} renderizado a HTML que incluya:

  \begin{itemize}
  \tightlist
  \item
    Limpieza documentada (faltantes, tipos, verificación de llaves).
  \item
    Tabla de KPIs por canal y región (ingreso, unidades, ticket
    promedio, participación).
  \item
    Al menos 3 gráficos con títulos que afirmen la conclusión.
  \item
    Una prueba estadística o modelo con interpretación.
  \item
    Al menos una función propia en \passthrough{\lstinline!R/!} usada en
    el reporte.
  \item
    Conclusiones y 3 recomendaciones accionables.
  \end{itemize}
\item
  Presentación oral de 5 minutos.
\end{enumerate}

\subsection{Rúbrica global del proyecto}

\begin{longtable}[]{@{}
  >{\raggedright\arraybackslash}p{(\columnwidth - 8\tabcolsep) * \real{0.2000}}
  >{\raggedright\arraybackslash}p{(\columnwidth - 8\tabcolsep) * \real{0.2000}}
  >{\raggedright\arraybackslash}p{(\columnwidth - 8\tabcolsep) * \real{0.2000}}
  >{\raggedright\arraybackslash}p{(\columnwidth - 8\tabcolsep) * \real{0.2000}}
  >{\raggedright\arraybackslash}p{(\columnwidth - 8\tabcolsep) * \real{0.2000}}@{}}
\toprule\noalign{}
\begin{minipage}[b]{\linewidth}\raggedright
Criterio (peso)
\end{minipage} & \begin{minipage}[b]{\linewidth}\raggedright
4 --- Excelente
\end{minipage} & \begin{minipage}[b]{\linewidth}\raggedright
3 --- Bueno
\end{minipage} & \begin{minipage}[b]{\linewidth}\raggedright
2 --- Suficiente
\end{minipage} & \begin{minipage}[b]{\linewidth}\raggedright
1 --- Insuficiente
\end{minipage} \\
\midrule\noalign{}
\endhead
\bottomrule\noalign{}
\endlastfoot
Reproducibilidad (20 \%) & Renderiza sin errores en otra computadora;
rutas relativas; \passthrough{\lstinline!renv!} & Renderiza con ajustes
menores & Requiere varios ajustes & No renderiza \\
Manejo de datos (20 \%) & Limpieza completa y justificada; verificación
de llaves & Limpieza correcta, poco documentada & Limpieza parcial &
Datos sin revisar \\
Transformación y KPIs (20 \%) & KPIs correctos, código
\passthrough{\lstinline!dplyr!} idiomático & Correctos con código
redundante & Errores menores & Errores graves \\
Visualización (15 \%) & Gráficos adecuados, claros y con mensaje &
Adecuados, mejorables & Tipo de gráfico poco apropiado & Ilegibles \\
Estadística (15 \%) & Método pertinente, supuestos revisados,
interpretación correcta & Pertinente, interpretación parcial & Método
dudoso & Ausente o incorrecto \\
Comunicación (10 \%) & Conclusiones y recomendaciones claras y
accionables & Claras pero generales & Confusas & Ausentes \\
\end{longtable}

\section{Referencias verificadas}

\begin{itemize}
\tightlist
\item
  \textbf{Web:} Posit / Quarto. \emph{Tutorial: Computations} (RStudio).
  \url{https://quarto.org/docs/get-started/computations/rstudio}
  {[}EN{]} {[}gratuito{]}
\item
  \textbf{Web:} Posit / Quarto. \emph{Using R}.
  \url{https://quarto.org/docs/computations/r.html} {[}EN{]}
  {[}gratuito{]}
\item
  \textbf{Web:} Posit / Quarto. \emph{Tutorial: Hello, Quarto}.
  \url{https://quarto.org/docs/get-started/hello/rstudio.html} {[}EN{]}
  {[}gratuito{]}
\item
  \textbf{Libro:} Wickham, H., Çetinkaya-Rundel, M., \& Grolemund, G.
  (2023). \emph{R for Data Science} (2.ª ed.). O'Reilly Media. ISBN
  978-1-4920-9740-2. Español: \url{https://davidrsch.github.io/r4ds-es/}
  {[}ES{]} {[}OA{]} --- caps. ``Quarto'' y ``Formatos de Quarto''.
\item
  \textbf{Artículo:} Wickham, H. (2014). Tidy data. \emph{Journal of
  Statistical Software, 59}(10), 1--23.
  \url{https://doi.org/10.18637/jss.v059.i10} {[}EN{]} {[}OA{]}
\end{itemize}
